\Askhsh{22150}{4}
\begin{ylh}{1}
    \item 8.2. Κριτήρια ομοιότητας
    \item 10.3. Εμβαδόν βασικών ευθύγραμμων σχημάτων
    \item 10.5. Λόγος εμβαδών όμοιων τριγώνων - πολυγώνων.
\end{ylh}
\begin{wrapfigure}[12]{r}{6.5cm} % Προσαρμόστε τον αριθμό γραμμών και το πλάτος ανάλογα με το κείμενο.
\centering
\begin{tikzpicture}[scale=0.9]
    % Ορισμός κορυφών
    \tkzDefPoint(0,4){A}
    \tkzDefPoint(-3,0){C} % Γ
    \tkzDefPoint(3,0){B}

    % Ορισμός μέσων
    \tkzDefMidPoint(C,B) \tkzGetPoint{D} % Δ
    \tkzDefMidPoint(A,C) \tkzGetPoint{E}
    \tkzDefMidPoint(A,B) \tkzGetPoint{Z}

    % Σχεδίαση πλευρών τριγώνου και εσωτερικών τμημάτων
    \tkzDrawPolygon(A,B,C) % Κύριο τρίγωνο ΑΒΓ
    \tkzDrawSegments(A,D E,Z E,D Z,D) % Εσωτερικά τμήματα

    % Σχεδίαση σημείων
    \tkzDrawPoints(A,B,C,D,E,Z)

    % Ετικέτες σημείων
    \tkzLabelPoint[above](A){A}
    \tkzLabelPoint[right](B){B}
    \tkzLabelPoint[left](C){$\Gamma$}
    \tkzLabelPoint[below](D){$\Delta$}
    \tkzLabelPoint[left](E){E}
    \tkzLabelPoint[right](Z){Z}
\end{tikzpicture}
\end{wrapfigure}
Δίνεται τρίγωνο ΑΒΓ με Ε και Ζ τα μέσα των πλευρών του ΑΓ και ΑΒ, αντίστοιχα.
\begin{alist}
\item Αν επιπλέον το ευθύγραμμο τμήμα ΑΔ ενώνει την κορυφή Α του τριγώνου ΑΒΓ και το μέσο Δ της απέναντι πλευράς ΒΓ, όπως στο σχήμα, να αποδείξετε ότι:
    \begin{rlist}
        \item Τα τρίγωνα ΕΔΓ και ΑΒΓ είναι όμοια με λόγο ομοιότητας $\dfrac{1}{2}$.
        \item Για το εμβαδόν $(ΑΕΔΖ)$ του τετραπλεύρου ΑΕΔΖ ισχύει ότι $(ΑΕΔΖ) = (ΑΒΓ) - 2(ΕΔΓ)$.
        \item Το εμβαδόν του τετραπλεύρου ΑΕΔΖ είναι ίσο με το $\dfrac{1}{2}$ του εμβαδού του τριγώνου ΑΒΓ. \monades{18}
    \end{rlist}
\item Αν το σημείο Δ είναι τυχαίο εσωτερικό σημείο της πλευράς ΒΓ του τριγώνου ΑΒΓ, τότε ισχύει ότι το εμβαδόν του τετραπλεύρου ΑΕΔΖ είναι ίσο με το $\dfrac{1}{2}$ του εμβαδού του τριγώνου ΑΒΓ; \monades{07}
\end{alist}